\documentclass[10pt]{article}
\usepackage[margin=25mm]{geometry}
\usepackage{amsmath,amssymb,hyperref}
\title{Implementation Notes for Federated LPV Identification}
\author{}
\date{}
\begin{document}
\maketitle
These notes accompany the manuscript and record numerical settings omitted
from its main explanation. They do not change the estimators, simulated
trajectories, search budgets, or evaluation results.

\section{Local output-error estimation}
The 44-s identification maneuver uses a four-second cosine ramp within the
initial 20-s settling interval, a four-second hold, and a final 20-s
identification segment. At $T_s=0.01$ s, only the final segment supplies the
2000 transitions used by OE. Both training and validation use noisy full-state
measurements, while identification uses known applied steering.

The affine model estimates the four entries of $\boldsymbol A$, two entries
of $\boldsymbol B$, the two-component offset, and the initial state.
Trust-region least squares uses analytic sensitivities and at most 300
function evaluations per start. Four initializations are compared: an affine
OLS estimate with initial spectral radius capped at 0.995,
$\boldsymbol A=0.9\boldsymbol I_2$, $\boldsymbol A=0.98\boldsymbol I_2$,
and the nominal 20-m/s model. The cap applies only to initialization.
No stability projection is imposed on the final OE estimate. The converged
solution with the lowest training loss is selected.

Trajectory checks require steering amplitude and rate bounds,
$|\beta|\leq5^\circ$, and $|a_y|\leq6$ m/s$^2$. During the final 2 s before
identification, the absolute difference between mean yaw rate and its reference must not exceed
$\max\{10^{-3}\,\mathrm{rad/s},0.01v\kappa_0\}$.
The absolute fitted slopes of $\beta/(0.25^\circ)$ and
$r/(2^\circ/\mathrm{s})$ must not exceed $0.02\,\mathrm{s}^{-1}$.
Limiter activation is allowed because the actual applied input is recorded.
A failed local fit or trajectory check invalidates the fleet condition;
no individual vehicles are removed or replaced.

\section{Basis normalization and smoothness penalty}
The raw basis is
\[
\boldsymbol b(v)=[1,v^{-1},v^{-2},z,z^2,z^3,z^4]^\top,
\qquad z=(v-20)/10.
\]
Stacking its rows on 401 equally spaced speeds from 10 to 30 m/s gives
$\boldsymbol B_{\rm ref}=\boldsymbol Q\boldsymbol R$ by reduced QR,
where $\boldsymbol Q\in\mathbb R^{401\times7}$ has orthonormal columns.
The common basis satisfies
$\boldsymbol\phi(v)^\top=\boldsymbol b(v)^\top\boldsymbol R^{-1}$.
The 401 nodes are numerical normalization and quadrature points, not
additional identification observations. Normalization preserves the function
span, but coefficient ridge penalties and reported condition numbers refer
to these normalized coordinates.

For Smooth Local, $\boldsymbol Q_2\in\mathbb R^{401\times7}$ factors the
trapezoidal approximation of
\[
\frac12\int_{-1}^{1}
\left\|\frac{d^2}{dz^2}
  [\boldsymbol\phi(20+10z)^\top\boldsymbol\Theta]\right\|_2^2\,dz.
\]
Affine functions in $z$ are unpenalized. The objective combines mean local
squared error, this curvature penalty with weight $\alpha$, and numerical
ridge $\eta=10^{-16}$. Development compared
$\alpha\in\{10^{-6},10^{-4},10^{-2},1,10^2,10^4,10^6\}$ and selected the
largest penalty within one standard error of the best fleet mean,
$\alpha=10^6$. Evaluation retains that setting, with all seven penalties
reported separately as sensitivity checks.

\section{Group regression and numerical equivalence}
The group search uses $K=1,\ldots,5$, one start for $K=1$, twelve seeded
K-means starts otherwise, at most 30 assignment iterations, and at least
five vehicles per group. Unchanged assignments stop iteration early.
Ties choose the lowest group index. A final changed assignment is followed
by a regression refit and score calculation without another reassignment.
The ridge term is $\rho=10^{-8}+10^{-12}$, including the numerical diagonal
adjustment. Assignment scores use entry scaling; regression does not.
Consequently, the procedure has no common-objective descent guarantee.

Assignment scales use the population standard deviations $\tau_j$ of target
entries over all observed vehicle--speed rows. Define
$s_j=\max\{\tau_j,0.05\operatorname{median}(\tau_1,\ldots,\tau_6)\}$
and $\boldsymbol D=\operatorname{diag}(s_1,\ldots,s_6)$.
The server computes moments using the population divisor $\sum_i n_i$,
checks all scales are positive, and fixes them for the search. K-means
initialization instead standardizes each mean-residual entry across clients
to zero mean and unit variance, using unit scale for constant entries.

For fixed group assignments, a separate centralized augmented least-squares
solve uses exactly the same targets, basis, and ridge. Coefficients and
speed-grid predictions must agree with the summed-statistics solution
within $10^{-8}$. This tests regression equivalence for the same partition,
not equality of independently selected groupings. The prespecified primary
comparison also requires all ten fleets to pass fit and trajectory checks,
and both adjusted confidence-interval lower bounds to exceed zero.

If $\boldsymbol\Phi_g$ stacks a group's design matrices, unregularized
regression has a unique minimizer precisely when
$\operatorname{rank}\boldsymbol\Phi_g=L$. This follows because
$\boldsymbol\Phi_g^\top\boldsymbol\Phi_g$ is positive definite exactly
under full column rank. Otherwise a nonzero coefficient perturbation in
the nullspace leaves all observed predictions unchanged. With exact targets
from a common coefficient matrix, full rank also gives exact recovery of
that matrix. Neither this conditional statement nor ridge regularization
proves that the assignment procedure finds informative groups.

\section{Control and reproducibility}
Controller synthesis uses the learned dynamic matrices and the common LQI
weights. OE offsets and initial states are not scheduled. Control evaluation
uses noise-free state feedback, a 12-degree steering bound, a 120-degree/s
rate bound, and a 6-degree sideslip bound. The identification controller's
rate bound is instead 30 degrees/s.
The complete variable-speed protocol and six-method results are recorded in
\texttt{docs/oe\_lqi\_validation.html} and
\texttt{report/oe\_lqi\_validation\_analysis.html}.
The independent evaluation configuration and local-estimation code are in
\texttt{code/config/paper\_oe\_lpv\_confirmation.json} and
\texttt{code/src/federated\_lpv/fixed\_lti\_identification.py}.
\end{document}
